\documentclass[11pt]{article}
\usepackage[margin=1in]{geometry}
\usepackage{amsmath,amssymb,bm}
\usepackage[T1]{fontenc}
\usepackage[utf8]{inputenc}
\usepackage{lmodern}
\usepackage{booktabs}
\usepackage{longtable}
\usepackage{array}
\usepackage{enumitem}
\usepackage{hyperref}

\title{Mount-In-Propagation Formulation Revision}
\author{sensor-fusion documentation}
\date{\today}

\newcommand{\R}{\mathbb{R}}
\newcommand{\T}{^{T}}
\newcommand{\dt}{\Delta t}
\newcommand{\crossmat}[1]{\left[#1\right]_\times}
\newcommand{\dq}{\delta q}

\begin{document}
\maketitle
\tableofcontents
\newpage

\section{Scope}

This document records the formulation change that moved mount alignment into
the inertial propagation model for both the runtime ESKF and the loose
reference filter. It is intentionally a before/after document. The full
steady-state formulations remain in the dedicated ESKF and loose formulation
PDFs, while this document explains why the revision was needed, what changed
mathematically, and how the implementation maps to the revised equations.

The motivating failure was observed on a synthetic figure-eight profile with
explicit vehicle roll excitation. With truth-quality inputs, both filters
previously converged to a near-correct composed vehicle attitude while retaining
a persistent mount roll error. The old formulation allowed vehicle roll error
and residual mount roll error to trade against each other. The revised
formulation removes that practical gauge by making mount errors affect
strapdown attitude and velocity propagation.

\section{Notation and Frames}

\begin{longtable}{>{\raggedright\arraybackslash}p{0.16\linewidth} >{\raggedright\arraybackslash}p{0.74\linewidth}}
\toprule
\textbf{Symbol} & \textbf{Meaning} \\
\midrule
\(b\) & Raw IMU body frame. \\
\(v\) & Physical vehicle frame, forward-right-down. \\
\(n\) & Local NED navigation frame used by the ESKF. \\
\(e\) & ECEF frame used by the loose reference filter. \\
\(s\) & Old seeded frame produced by pre-rotating raw IMU with a coarse mount seed. \\
\(c\) & Old corrected vehicle frame after applying residual mount. \\
\(q_{vb}\) & Physical vehicle-to-body mount quaternion, with \(x_b=C_{bv}x_v\). \\
\(q_{nv}\) & Vehicle-to-NED attitude quaternion, with \(x_n=C_{nv}x_v\). \\
\(q_{ev}\) & Vehicle-to-ECEF attitude quaternion, with \(x_e=C_{ev}x_v\). \\
\bottomrule
\end{longtable}

Rotations follow the project convention
\begin{equation}
  x_a = C_{ab}x_b,
  \qquad
  C(q_1\otimes q_2)=C(q_1)C(q_2).
\end{equation}

Small-angle perturbations use the first-order quaternion
\begin{equation}
  \dq(\delta\theta) \approx
  \begin{bmatrix}1 & \frac{1}{2}\delta\theta\T\end{bmatrix}\T .
\end{equation}

\section{Old Formulation}

\subsection{State Split}

The old ESKF nominal state was
\begin{equation}
  x_{\mathrm{old}} =
  \begin{bmatrix}
  q_{ns} & v_n\T & p_n\T & b_g\T & b_a\T & q_{cs}
  \end{bmatrix}\T ,
\end{equation}
where \(q_{ns}\) rotated the seeded frame into NED and \(q_{cs}\) rotated the
seeded frame into the corrected vehicle frame. The corresponding error state
was
\begin{equation}
  \delta x_{\mathrm{old}} =
  \begin{bmatrix}
  \delta\theta_s &
  \delta v_n &
  \delta p_n &
  \delta b_g &
  \delta b_a &
  \delta\psi_{cs}
  \end{bmatrix}\T .
\end{equation}

The loose filter had the same structure in ECEF:
\begin{equation}
  x_{\mathrm{old,loose}} =
  \begin{bmatrix}
  q_{es} & v_e\T & p_e\T & b_g\T & b_a\T &
  s_g\T & s_a\T & q_{cs}
  \end{bmatrix}\T .
\end{equation}

\subsection{IMU Pre-Rotation}

The align filter provided a coarse physical mount seed \(q_{vb}^{0}\). Before
the ESKF or loose propagation step, the runtime converted raw IMU samples into
the seeded frame:
\begin{align}
  \Delta\theta_s &= C_{bv}^{0\T}\Delta\theta_b,\\
  \Delta v_s &= C_{bv}^{0\T}\Delta v_b.
\end{align}
The filters then propagated using \(\Delta\theta_s\) and \(\Delta v_s\). The
residual mount \(q_{cs}\) was held constant during propagation:
\begin{equation}
  q_{cs}^{+}=q_{cs}^{-}.
\end{equation}

\subsection{Old ESKF Propagation}

The old ESKF prediction was effectively
\begin{align}
  \Delta\theta_s^c &= \Delta\theta_s - b_g\dt,\\
  \Delta v_s^c &= \Delta v_s - b_a\dt,\\
  q_{ns}^{+} &= \operatorname{normalize}
    \left(q_{ns}^{-}\otimes\dq(\Delta\theta_s^c)\right),\\
  v_n^{+} &= v_n^{-}+C_{ns}\Delta v_s^c+g_n\dt,\\
  p_n^{+} &= p_n^{-}+v_n^{-}\dt.
\end{align}

The important structural point is that \(q_{cs}\) did not appear in these
equations. Therefore the generated propagation Jacobian had no dynamic
sensitivity from mount error into attitude or velocity:
\begin{equation}
  \frac{\partial q_{ns}^{+}}{\partial \delta\psi_{cs}} \approx 0,
  \qquad
  \frac{\partial v_n^{+}}{\partial \delta\psi_{cs}} \approx 0.
\end{equation}

\subsection{Old NHC Measurement}

The old non-holonomic constraint used the corrected vehicle velocity
\begin{equation}
  v_c = C_{cs}C_{ns}\T v_n .
\end{equation}
The lateral and vertical constraints were
\begin{equation}
  z_y = e_y\T v_c = 0,\qquad z_z = e_z\T v_c = 0.
\end{equation}
This gave a direct residual-mount measurement Jacobian:
\begin{equation}
  H_{\psi_{cs}} =
  \frac{\partial(e_i\T C_{cs}C_{ns}\T v_n)}
       {\partial \delta\psi_{cs}}.
\end{equation}

\subsection{Why This Was Insufficient}

The old model could produce a correct composed vehicle-frame quantity while
splitting the error incorrectly:
\begin{equation}
  C_{cv}=C_{cs}C_{sn}
\end{equation}
could remain nearly correct even if the attitude and residual mount separately
were wrong. In the observed synthetic roll-excitation case, the filters settled
to approximately
\begin{align}
  \text{mount roll error} &\approx +0.45^\circ,\\
  \text{vehicle roll error} &\approx -0.45^\circ,
\end{align}
with a near-zero sum. Since \(q_{cs}\) did not affect propagation, the IMU
dynamics did not provide an independent signal to decide whether the correction
belonged in attitude or mount. Tuning covariance could change allocation, but
it could not remove the structural ambiguity.

\section{Revised Formulation}

\subsection{State Meaning}

The revised ESKF keeps the same state dimension and Rust field layout, but
changes the meaning of the quaternion formerly named \texttt{qcs}:
\begin{equation}
  q_{cs}\quad\hbox{is now interpreted as}\quad q_{vb}.
\end{equation}
That is, \texttt{qcs0..qcs3} store the physical vehicle-to-body mount, not a
residual seeded-frame-to-vehicle correction.

The revised ESKF nominal state is
\begin{equation}
  x_{\mathrm{new}} =
  \begin{bmatrix}
  q_{nv} & v_n\T & p_n\T & b_g\T & b_a\T & q_{vb}
  \end{bmatrix}\T ,
\end{equation}
and the revised error state is
\begin{equation}
  \delta x_{\mathrm{new}} =
  \begin{bmatrix}
  \delta\theta_v &
  \delta v_n &
  \delta p_n &
  \delta b_g &
  \delta b_a &
  \delta\psi_{vb}
  \end{bmatrix}\T .
\end{equation}

The loose filter analog is
\begin{equation}
  x_{\mathrm{new,loose}} =
  \begin{bmatrix}
  q_{ev} & v_e\T & p_e\T & b_g\T & b_a\T &
  s_g\T & s_a\T & q_{vb}
  \end{bmatrix}\T .
\end{equation}

\subsection{Raw IMU Propagation Through Mount}

The filters now consume raw body-frame IMU increments. Bias-corrected body-frame
increments are
\begin{align}
  \Delta\theta_b^c &= \Delta\theta_b - b_g\dt,\\
  \Delta v_b^c &= \Delta v_b - b_a\dt.
\end{align}
The current mount rotates them into the vehicle frame:
\begin{align}
  \Delta\theta_v^c &= C_{bv}(q_{vb})\T \Delta\theta_b^c,\\
  \Delta v_v^c &= C_{bv}(q_{vb})\T \Delta v_b^c.
\end{align}

The revised ESKF propagation is
\begin{align}
  q_{nv}^{+} &=
    \operatorname{normalize}\left(q_{nv}^{-}\otimes
    \dq(\Delta\theta_v^c)\right),\\
  v_n^{+} &= v_n^{-}+C_{nv}\Delta v_v^c+g_n\dt,\\
  p_n^{+} &= p_n^{-}+v_n^{-}\dt,\\
  q_{vb}^{+} &= q_{vb}^{-}.
\end{align}

For loose, the same structure is used with ECEF gravity and Earth rotation
terms:
\begin{align}
  \omega_v &= C_{bv}(q_{vb})\T\omega_b^c,\\
  f_v &= C_{bv}(q_{vb})\T f_b^c,\\
  \dot q_{ev} &=
    \frac{1}{2}q_{ev}\otimes
    \begin{bmatrix}0 & \omega_v\T\end{bmatrix}\T
    + \hbox{Earth-rate correction},\\
  \dot v_e &= C_{ev}f_v + g_e(p_e) - 2\omega_{ie,e}\times v_e,\\
  \dot p_e &= v_e.
\end{align}

\subsection{New Propagation Sensitivities}

The essential new Jacobian blocks are the mount sensitivities in propagation.
For a small mount perturbation,
\begin{equation}
  C_{bv}(\dq(\delta\psi_{vb})\otimes q_{vb})\T x_b
  \approx
  C_{bv}(q_{vb})\T x_b
  - C_{bv}(q_{vb})\T \crossmat{x_b}\delta\psi_{vb}.
\end{equation}
Therefore, to first order,
\begin{align}
  \frac{\partial \Delta\theta_v^c}{\partial \delta\psi_{vb}}
    &\approx -C_{bv}\T\crossmat{\Delta\theta_b^c},\\
  \frac{\partial \Delta v_v^c}{\partial \delta\psi_{vb}}
    &\approx -C_{bv}\T\crossmat{\Delta v_b^c}.
\end{align}

These terms create covariance between mount error, attitude error, and velocity
error during propagation. GNSS position/velocity and NHC can then correct mount
through this propagated cross-covariance. This is the observability path that
was missing before.

\subsection{Revised NHC Measurement}

Since \(q_{nv}\) now represents the physical vehicle attitude, NHC no longer
needs the mount quaternion directly:
\begin{equation}
  v_v = C_{nv}\T v_n.
\end{equation}
The ESKF NHC rows are
\begin{equation}
  z_y=e_y\T C_{nv}\T v_n=0,\qquad
  z_z=e_z\T C_{nv}\T v_n=0.
\end{equation}
The loose rows are
\begin{equation}
  z_y=e_y\T C_{ev}\T v_e=0,\qquad
  z_z=e_z\T C_{ev}\T v_e=0.
\end{equation}
Thus
\begin{equation}
  \frac{\partial z_i}{\partial\delta\psi_{vb}} = 0
\end{equation}
for the instantaneous NHC row. This is expected and desired. Mount is no
longer directly nudged by NHC in a single update; it is corrected through the
state correlations generated by propagation.

\section{Before/After Comparison}

\begin{longtable}{>{\raggedright\arraybackslash}p{0.28\linewidth} >{\raggedright\arraybackslash}p{0.31\linewidth} >{\raggedright\arraybackslash}p{0.31\linewidth}}
\toprule
\textbf{Aspect} & \textbf{Before} & \textbf{After} \\
\midrule
Mount state meaning &
Residual seeded-frame-to-vehicle \(q_{cs}\). &
Physical vehicle-to-body \(q_{vb}\), stored in the same Rust fields. \\
\addlinespace
IMU input to propagation &
Raw body IMU was pre-rotated by align seed before entering filter. &
Raw body IMU enters filter; generated propagation rotates through current \(q_{vb}\). \\
\addlinespace
Mount in propagation &
No direct appearance in attitude or velocity propagation. &
Directly affects attitude and velocity propagation through \(C_{bv}\T\). \\
\addlinespace
NHC measurement &
\(v_c=C_{cs}C_{ns}\T v_n\), direct residual-mount Jacobian. &
\(v_v=C_{nv}\T v_n\), no direct mount Jacobian. \\
\addlinespace
Mount observability path &
Mostly NHC direct row plus covariance artifacts. &
Propagation cross-covariance plus GNSS/NHC corrections. \\
\addlinespace
Failure mode &
Vehicle attitude and residual mount could absorb equal/opposite errors. &
Wrong mount causes wrong inertial propagation, so residuals push mount. \\
\bottomrule
\end{longtable}

\section{Implementation Mapping}

\subsection{Generated ESKF}

The symbolic source is \path{ekf/eskf.py}. The generated Rust fragments are
included by \path{ekf/src/generated_eskf.rs}. The important changes are:
\begin{itemize}[leftmargin=*]
  \item \texttt{propagate\_nominal(...)} now takes \texttt{qcs} as \(q_{vb}\).
  \item Nominal prediction rotates \(\Delta\theta_b\) and \(\Delta v_b\)
        through \(C_{bv}\T\).
  \item Error transition and noise input include mount-to-attitude and
        mount-to-velocity propagation sensitivity.
  \item Body velocity observations no longer depend on \texttt{qcs}.
\end{itemize}

\subsection{Runtime ESKF Facade}

\texttt{ekf/src/fusion.rs} no longer pre-rotates IMU before calling
\texttt{RustEskf::predict}. It initializes \texttt{nominal.qcs*} from the align
or reference mount seed. In follow-align mode, the facade writes the current
align mount into \texttt{nominal.qcs*} and keeps mount error states frozen.

\subsection{Loose Reference Filter}

The symbolic source is \texttt{ekf/ins\_gnss\_loose.py}. The generated Rust
fragments are included by \texttt{ekf/src/generated\_loose.rs}. The important
changes are:
\begin{itemize}[leftmargin=*]
  \item The transition matrix uses \texttt{qcs} as \(q_{vb}\) in the IMU
        rotation terms.
  \item Loose nominal propagation in \texttt{ekf/src/loose.rs} rotates raw
        body-frame gyro and accelerometer measurements into the vehicle frame.
  \item Loose NHC rows observe \(C_{ev}\T v_e\) and have no direct mount
        support.
  \item The visualizer seeds loose \texttt{qcs} with align mount and plots it
        directly as physical mount.
\end{itemize}

\section{Covariance and Initialization Implications}

Moving mount into propagation makes mount observable through cross-covariance,
so the initial covariance must allow the filter to adjust the attitude/mount
split. This matters most for the loose filter because it initializes vehicle
attitude from a yaw-only state. If the loose vehicle-attitude covariance is too
small, the filter remains overconfident in an attitude that cannot yet be known
well enough, and the mount state cannot fully absorb the coarse seed error.

The visualizer default was therefore changed from
\begin{equation}
  \sigma_{\theta,\mathrm{loose}}=2^\circ
\end{equation}
to
\begin{equation}
  \sigma_{\theta,\mathrm{loose}}=20^\circ.
\end{equation}
This is not an ad-hoc correction to one maneuver. It reflects the weaker
initial information content of the loose yaw-only attitude seed.

\section{Synthetic Validation}

On the truth-input figure-eight roll-excitation profile without GNSS fault, the
old filters settled at nonzero mount roll error:
\begin{center}
\begin{tabular}{lcc}
\toprule
Filter & Old final mount qerr & New final mount qerr \\
\midrule
ESKF & approximately \(0.46^\circ\) & \(0.0277^\circ\) \\
Loose & approximately \(0.45^\circ\) & \(0.0321^\circ\) \\
\bottomrule
\end{tabular}
\end{center}

With an early GNSS velocity fault of \(+0.5\,\mathrm{m/s}\) north over
\(0\ldots360\,\mathrm{s}\), the revised formulation still converges:
\begin{center}
\begin{tabular}{lc}
\toprule
Filter & Final mount qerr \\
\midrule
ESKF & \(0.0887^\circ\) \\
Loose & \(0.0545^\circ\) \\
\bottomrule
\end{tabular}
\end{center}
In that same faulted case, the loose GNSS position gate accepted all GNSS
position updates after the revision. Before this change, the bad state split
could drive the loose prediction far enough away that the full 3D position gate
rejected valid measurements for the rest of the run.

\section{Test Coverage}

The following tests exercise the revised model:
\begin{itemize}[leftmargin=*]
  \item \texttt{ekf/tests/eskf\_nhc\_jacobian.rs}: finite-difference check for
        the direct vehicle-frame ESKF NHC Jacobians.
  \item \texttt{ekf/tests/nhc\_eskf\_loose\_equivalence.rs}: verifies ESKF and
        loose NHC rows and updates remain equivalent after the attitude basis
        transform.
  \item \texttt{ekf/tests/loose\_mount\_observability.rs}: verifies NHC does
        not directly move mount under the new formulation; mount is corrected
        through propagation correlations.
  \item The synthetic GNSS/INS test verifies that roll-excitation now makes
        internal mount estimates converge.
\end{itemize}

\section{Remaining Follow-Up}

The existing ESKF and loose formulation PDFs still describe portions of the old
seeded-frame terminology. They should be updated to use this document as the
source of truth for mount semantics, or merged into a single current
formulation reference after the revised behavior is accepted on the full
experimental data suite.

\end{document}
